\rtmtight
\begin{tblr}{width=\textwidth,colspec={|Q[c,m,2.1cm]|X[l,m]|},hlines,vlines,hspan=minimal,rowsep=1.5pt,colsep=3pt,row{1}={bg=headgray,font=\bfseries}}
\SetCell{c,m}\hfil\logo\hfil & \textbf{POLITEKNIK NEGERI MALANG}\par\textbf{JURUSAN TEKNOLOGI INFORMASI}\par\textbf{PROGRAM STUDI: D4 TEKNIK INFORMATIKA} \\
\SetCell[c=2]{l} \textbf{RENCANA TUGAS MAHASISWA (RTM) DAN RUBRIK PENILAIAN} & \\
Nama Mata Kuliah & Agama \\
Kode Mata Kuliah & RTI252001 \quad \textbf{sks / Jam perminggu:} 2 SKS/2 jam \quad \textbf{Semester:} 2 \\
Dosen Pengampu & \begin{enumerate}\item Astrifidha Rahma Amalia, S.Pd., M.Pd.\end{enumerate} \\
\end{tblr}
\assessmentblock{Case Method Diskusi dan Presentasi Kelompok}{Case Method (lembar observasi dosen)}{SCPMK0102-01001, SCPMK0102-01002, SCPMK0102-01004}{Mahasiswa mendiskusikan kasus keagamaan aktual pada topik manusia dan agama, tauhid, ekologi, akhlak, IPTEKS, etos kerja, ekonomi syariah, dan masyarakat Islam, lalu mempresentasikan hasil kelompok.}{Berdiskusi secara kelompok, menyusun makalah, mempresentasikan lewat PPT, melakukan tanya jawab, dan menghubungkan materi ajar dengan kehidupan sehari-hari.}{Makalah kelompok, bahan presentasi PPT, dan catatan hasil diskusi.}{Partisipasi aktif dalam diskusi kelompok & 15 & Hadir penuh, mengajukan pendapat konstruktif di setiap sesi diskusi, serta merespons pendapat anggota lain dengan argumen. & Ikut berdiskusi, tetapi pendapat jarang disampaikan atau respons terhadap anggota lain terbatas. & Pasif, tidak memberikan pendapat, atau tidak terlibat dalam kerja kelompok. \\ \\
Ketepatan menghubungkan konsep dengan dalil dan contoh & 15 & Konsep manusia, agama, tauhid, hingga masyarakat Islam dihubungkan dengan dalil Al-Qur'an/Hadis dan contoh kehidupan nyata yang relevan. & Konsep dijelaskan benar, tetapi dalil atau contoh yang diberikan kurang relevan atau tidak lengkap. & Penjelasan konsep keliru atau tidak disertai dalil dan contoh sama sekali. \\ \\
Keterbukaan dan etika dalam menanggapi pendapat & 10 & Mendengarkan, mempertimbangkan, dan menanggapi pandangan berbeda dengan sopan, hormat, dan menghargai perbedaan keyakinan. & Menanggapi pendapat lain dengan cukup sopan, tetapi kadang memotong atau mengabaikan pandangan berbeda. & Menolak pandangan lain tanpa argumen, bersikap tidak sopan, atau merendahkan keyakinan berbeda. \\ \\
Kejelasan penyampaian ide dan kualitas presentasi & 15 & Makalah dan PPT terstruktur, ide disampaikan jelas dan efektif, serta tanya jawab dijawab dengan komunikasi profesional. & Presentasi cukup jelas, tetapi struktur makalah/PPT atau jawaban tanya jawab masih kurang runtut. & Presentasi sulit dipahami, makalah/PPT tidak terstruktur, atau tidak mampu menjawab pertanyaan. \\}

\assessmentblock{Tugas Resume Konsep Manusia}{Tugas (esai/resume)}{SCPMK0102-01002}{Mahasiswa menyusun resume mandiri tentang konsep manusia dalam pandangan Islam, fungsi dan tugas hidup manusia, serta hakikat kehidupan dunia dan akhirat.}{Menyaring informasi relevan dari berbagai sumber, memilih poin terpenting, menyusunnya dalam format terstruktur dan ringkas, serta berpikir kritis terhadap sumber yang digunakan.}{Resume tertulis (esai) yang dikumpulkan sesuai tenggat.}{Cakupan konsep utama dalam resume & 4 & Resume mencakup konsep manusia, fungsi dan tugas hidup, serta hakikat dunia-akhirat secara akurat sesuai materi dan sumber utama. & Sebagian besar konsep utama tercakup, tetapi ada bagian yang dangkal atau kurang akurat. & Resume tidak mencakup konsep utama atau banyak isi yang keliru. \\ \\
Struktur dan kejelasan bahasa & 3 & Resume terorganisasi rapi, ringkas, menggunakan bahasa yang baik, dan mudah diikuti alurnya. & Struktur cukup jelas, tetapi ada bagian yang bertele-tele atau bahasa kurang rapi. & Resume tidak terstruktur, sulit dipahami, atau sekadar salinan sumber. \\ \\
Keterkaitan dengan kehidupan dan orisinalitas & 3 & Materi dihubungkan dengan pengalaman atau kasus kehidupan sehari-hari dalam konteks moral, etika, atau spiritual dengan pemahaman baru yang orisinal. & Ada upaya menghubungkan materi dengan kehidupan, tetapi analisisnya terbatas atau kurang orisinal. & Tidak ada keterkaitan dengan kehidupan nyata dan resume hanya menyalin tanpa pemahaman. \\}

\assessmentblock{Kuis 1: Tes Lisan Bacaan dan Hafalan Al-Qur'an}{Kuis (tes lisan)}{SCPMK0102-01002}{Mahasiswa membaca dan menghafalkan ayat-ayat Al-Qur'an yang berkaitan dengan materi tauhid, lalu dinilai secara lisan.}{Berlatih membaca sesuai tajwid, menghafalkan ayat yang ditugaskan, menjaga fokus dan ketelitian, kemudian menyetorkan bacaan dan hafalan secara lisan di hadapan dosen.}{Setoran lisan bacaan dan hafalan ayat Al-Qur'an.}{Ketepatan tajwid dan kelancaran bacaan & 3 & Bacaan sesuai aturan tajwid (pengucapan huruf, panjang-pendek suara, hukum bacaan) dan lancar tanpa terbata-bata. & Bacaan sebagian besar benar, tetapi ada beberapa kesalahan tajwid atau bacaan kadang terhenti. & Banyak kesalahan tajwid dan bacaan tidak lancar. \\ \\
Ketepatan hafalan dan penghayatan & 2 & Hafalan ayat tepat dan konsisten, dilantunkan dengan penghayatan, serta hafalan sebelumnya tetap terjaga. & Hafalan sebagian besar benar, tetapi ada bagian yang terlupa atau penghayatan kurang. & Hafalan banyak yang salah atau tidak mampu melafalkan ayat yang ditugaskan. \\}

\assessmentblock{UTS Pemahaman Konsep Manusia, Agama, dan Tauhid}{UTS (ujian tulis, closed book)}{SCPMK0102-01001, SCPMK0102-01002}{Mahasiswa mengerjakan ujian tulis materi minggu 1--7: pendahuluan, manusia dan agama, konsep tauhid, dan konsep manusia.}{Mengerjakan soal ujian secara mandiri, jujur, dan closed book sesuai jadwal asesmen.}{Lembar jawaban ujian tulis.}{Ketepatan penjelasan konsep manusia, agama, dan tauhid & 5 & Jawaban menjelaskan konsep manusia, agama, dan tauhid secara benar dan lengkap sesuai materi kuliah. & Jawaban sebagian benar, tetapi ada konsep yang tertukar atau penjelasan tidak lengkap. & Jawaban keliru atau tidak menjawab konsep yang ditanyakan. \\ \\
Ketepatan penggunaan dalil Al-Qur'an dan Hadis & 3 & Dalil yang dikutip relevan dengan konsep yang dijelaskan dan penerapannya tepat. & Dalil disebutkan, tetapi relevansi atau penerapannya pada jawaban kurang tepat. & Tidak menggunakan dalil atau dalil yang dikutip tidak berkaitan dengan soal. \\ \\
Koherensi argumen dan kejelasan tulisan & 2 & Argumen runtut, konsisten antarbagian jawaban, dan tulisan mudah dipahami. & Argumen cukup jelas, tetapi ada bagian yang melompat atau tidak konsisten. & Jawaban tidak runtut, kontradiktif, atau sulit dipahami. \\}

\assessmentblock{Kuis 2: Video Simulasi Etos Kerja}{Kuis (rubrik penilaian video simulasi)}{SCPMK0102-01003}{Mahasiswa membuat video simulasi (cosplay) tentang etos kerja dalam Islam beserta dalil Al-Qur'an dan Hadis yang relevan, dikumpulkan tepat waktu.}{Memahami konsep etos kerja, merangkum informasi menjadi satu cerita/demonstrasi, menggunakan perangkat lunak animasi, pengeditan video, atau desain grafis, lalu mengumpulkan video sesuai tenggat.}{Video simulasi beserta dalil pendukung.}{Kesesuaian isi video dengan konsep etos kerja dan dalil & 3 & Video mencerminkan pemahaman mendalam tentang etos kerja dalam Islam, dikaitkan dengan dalil Al-Qur'an/Hadis yang tepat, dan pesan tersampaikan sesuai tujuan pembelajaran. & Isi video cukup sesuai topik, tetapi kaitan dalil atau kedalaman pemahaman masih terbatas. & Isi video tidak menggambarkan konsep etos kerja atau tanpa dalil pendukung. \\ \\
Kreativitas, teknologi, dan ketepatan durasi & 2 & Video mengandung elemen kreatif, memanfaatkan animasi/grafik/perangkat lunak pengeditan, durasi proporsional, dan dikumpulkan tepat waktu. & Video dibuat dengan teknologi sederhana, kreativitas atau durasi kurang proporsional, tetapi tetap dikumpulkan tepat waktu. & Video seadanya tanpa unsur kreatif, durasi tidak sesuai materi, atau terlambat dikumpulkan. \\}

\assessmentblock{UAS Refleksi Nilai Keagamaan}{UAS (ujian tulis, closed book)}{SCPMK0102-01004}{Mahasiswa mengerjakan ujian tulis materi minggu 9--16: wawasan ekologi, aktualisasi akhlak, IPTEKS, etos kerja, ekonomi syariah, dan masyarakat Islam.}{Mengerjakan soal ujian secara mandiri, jujur, dan closed book sesuai jadwal asesmen.}{Lembar jawaban ujian tulis.}{Pemahaman etos kerja, ekonomi syariah, dan masyarakat Islam & 6 & Jawaban menjelaskan etos kerja, prinsip dan etika ekonomi syariah, serta konsep keluarga dan masjid secara benar dan lengkap. & Sebagian konsep dijelaskan benar, tetapi ada bagian yang dangkal atau tertukar. & Jawaban keliru atau tidak menunjukkan pemahaman materi. \\ \\
Refleksi nilai keagamaan pada kasus keputusan etis & 5 & Nilai keagamaan diterapkan secara tepat untuk menganalisis kasus keputusan etis di lingkungan kampus dengan pertimbangan yang jelas. & Refleksi nilai pada kasus cukup relevan, tetapi analisis atau pertimbangan etisnya belum mendalam. & Tidak mampu mengaitkan nilai keagamaan dengan kasus atau jawaban tidak reflektif. \\ \\
Penggunaan dalil dan koherensi argumen & 4 & Dalil Al-Qur'an/Hadis dikutip relevan dan argumen disusun runtut serta konsisten. & Dalil atau argumen ada, tetapi relevansi atau keruntutannya kurang. & Tidak menggunakan dalil dan argumen tidak runtut. \\}

\begin{tblr}{width=\textwidth,colspec={|Q[l,m,3.2cm]|X[l,m]|},hlines,vlines,hspan=minimal,row{1-3}={bg=headgray},rowsep=1.5pt,colsep=3pt}
Jadwal Pelaksanaan: & Minggu 1--17 sesuai RPS. Setiap evaluasi memiliki RTM dan rubrik multi-kriteria. \\
Lain-Lain Yang Diperlukan: & LMS, modul PAI, Al-Qur'an dan terjemah, perangkat pembuatan video, dan perangkat presentasi. \\
Pustaka: & Mengacu pustaka utama dan pendukung pada bagian RPS. \\
\end{tblr}
